\documentclass[12pt]{article}
\usepackage[utf8]{inputenc}
\usepackage[T1]{fontenc}
\usepackage[margin=1in]{geometry}
\usepackage{setspace}
\usepackage{xcolor}
\usepackage[colorlinks=true, urlcolor=blue!60!black, linkcolor=blue!60!black]{hyperref}
\setlength{\parindent}{0em}
\setlength{\parskip}{0.45em}
\renewcommand{\baselinestretch}{1.0}
\begin{document}
\begin{center}
{\Large \textit{Virginia v. Black}, 538 U.S. 343 (2003)\par}
\vspace{0.3em}
{\normalsize Case Brief --- One-Page Edition\par}
\end{center}

\textbf{Memory Jogger.} The cross-burning case: intimidation is an unprotected true threat, but intent cannot be presumed from the act (\href{https://supreme.justia.com/cases/federal/us/538/343/}{538 U.S. 343}).

\textbf{Facts.} Virginia banned cross burning with intent to intimidate and made any burning prima facie evidence of intent. Barry Black burned a cross at a Klan rally; Elliott and O'Mara tried to burn one on a Black neighbor's lawn.

\textbf{Procedural History.} The Supreme Court of Virginia struck the statute down under \textit{R.A.V.}

\textbf{Votes.} O'Connor wrote for the Court and a four-Justice plurality; Black's conviction fell 7--2, and the others were vacated 5--4. Thomas dissented.

\textbf{Holding.} A State may ban cross burning with intent to intimidate; the plurality held the prima facie provision, as construed, facially invalid.

\textbf{Reasoning.} Intimidation is a true threat, and \textit{R.A.V.} allows singling out cross burning, a ``particularly virulent form of intimidation,'' when the ban applies regardless of viewpoint. Presuming intent ignores context such as political rallies. Souter saw viewpoint suppression; Thomas called cross burning ``conduct, not expression.''

\textbf{Examples.} \emph{Same-Side:} a ban on nooses hung to intimidate. \emph{Opposite-Side:} convicting a protest artist on the presumption alone. \emph{Fence-Sitter:} a burning-cross image tagging a candidate online.

\textbf{Critique.} A cross-specific ban may blur harm and viewpoint.

\textbf{Key Quotation.} ``\,`True threats' encompass those statements where the speaker means to communicate a serious expression of an intent to commit an act of unlawful violence to a particular individual or group of individuals. \dots Intimidation in the constitutionally proscribable sense of the word is a type of true threat, where a speaker directs a threat to a person or group of persons with the intent of placing the victim in fear of bodily harm or death.'' --- Justice O'Connor, Opinion of the Court, 538 U.S. at 359--360.

\end{document}
